\documentclass[UTF8,a4paper,11pt,openany]{ctexbook}
\usepackage{../common/statlearnbook}
\SLSetBookMetadata{第 5 册}{统计学习方法讲义 v2.7.0}{采样方法、主题模型与图排序}
\begin{document}
\setcounter{page}{666}
\setcounter{chapter}{33}
\setcounter{figure}{2}
\pagestyle{headings}
\chaptermark{Gibbs 抽样}
\sectionmark{逐变量 Gibbs 抽样}
\textbf{学习算法。}在系统扫描中，第$t$轮从$x^{(t-1)}$开始，先抽$x_1^{(t)}$，再抽$x_2^{(t)}$；第$j$步使用本轮已经更新的$x_1^{(t)},\ldots,x_{j-1}^{(t)}$与尚未更新的$x_{j+1}^{(t-1)},\ldots,x_d^{(t-1)}$。图\ref{fig:V5-C04-coordinate-sweep}明确区分“轮内状态”和“轮末样本”。算法及其前置说明另行区分完整轮次与中途停止状态。
% v2.7.0 figure UID: FIG-P634-01
% Structured glyph-safe redraw: preserve the one-arrow/eight-slot semantics,
% but express low-profile punctuation and equality through grouping and arrows.
\tikzset{
  slfig-FIG-P634-01/.style={font=\fontsize{9.6pt}{11.5pt}\selectfont,every path/.append style={line cap=round,line join=round}},
  sl634-old/.style={draw=SLRuleGray,densely dotted,fill=white,minimum width=13.2mm,minimum height=9.5mm,line width=.78pt,inner sep=2pt},
  sl634-done/.style={draw=SLBlue,fill=SLBlue!4,pattern=north east lines,pattern color=SLRuleGray,minimum width=13.2mm,minimum height=9.5mm,line width=.95pt,inner sep=2pt},
  sl634-halo/.style={draw=none,fill=white,inner xsep=1.5pt,inner ysep=1.25pt},
  sl634-now/.style={draw=SLGold,fill=SLGold!9,minimum width=13.2mm,minimum height=9.5mm,line width=1.05pt,inner sep=2pt},
  sl634-card/.style={draw=SLRuleGray,rounded corners=1.8pt,fill=white,line width=.55pt,inner xsep=5pt}
}
\pgfkeysifdefined{/tikz/alt}{}{\tikzset{alt/.code={}}}
\begin{figure}[htbp]
\centering
\begin{tikzpicture}[slfig-FIG-P634-01,>=Stealth,every node/.style={font=\fontsize{9.6pt}{11.5pt}\selectfont,align=center},
  alt={对象—关系—结论：固定顺序为 1,2,…,j−1,j,j+1,…,d；当前第 j 步的 x^[j] 前段是本轮新值且后段是前轮旧值；x^[d]=x^(t) 仅在末位更新结束后成立，并记录为轮末样本}]
  \node[font=\fontsize{10.6pt}{12.7pt}\selectfont\bfseries,text=SLBlue] at (0,2.16) {单轮系统扫描坐标带};
  \node at (-5.2,1.48) {$1$};
  \node at (-3.7,1.48) {$2$};
  \node at (-2.2,1.48) {省略};
  \node at (-.7,1.48) {前位};
  \node at (.8,1.48) {当前};
  \node at (2.3,1.48) {后位};
  \node at (3.8,1.48) {省略};
  \node at (5.3,1.48) {末位};
  \draw[-{Stealth[length=1.8mm,width=1.25mm]},draw=SLGray,line width=.60pt] (-5.52,1.10)--(5.62,1.10)
    node[right=4pt,font=\fontsize{9.6pt}{11.5pt}\selectfont] {更新顺序};
  \node[sl634-done] at (-5.2,0) {};
  \node[sl634-done] at (-3.7,0) {};
  \node[sl634-done] at (-2.2,0) {};
  \node[sl634-done] at (-.7,0) {};
  \node[sl634-halo] at (-5.2,0) {坐标\\$1$};
  \node[sl634-halo] at (-3.7,0) {坐标\\$2$};
  \node[sl634-halo] at (-2.2,0) {中间\\坐标};
  \node[sl634-halo] at (-.7,0) {前段\\末位};
  \node[sl634-now] at (.8,0) {当前\\坐标};
  \node[sl634-old] at (2.3,0) {后段\\首位};
  \node[sl634-old] at (3.8,0) {中间\\坐标};
  \node[sl634-old] at (5.3,0) {末位\\坐标};
  \node[font=\fontsize{9.6pt}{11.5pt}\selectfont,text=SLBlue] at (-2.9,-.76) {本轮新值};
  \node[font=\fontsize{9.6pt}{11.5pt}\selectfont,text=SLGold] at (.8,-.76) {当前新值};
  \node[font=\fontsize{9.6pt}{11.5pt}\selectfont,text=SLTextGray] at (3.8,-.76) {前轮旧值};

  \node[sl634-card,minimum width=118mm,minimum height=10.5mm] at (0,-1.52) {};
  \node[font=\fontsize{10.0pt}{12.0pt}\selectfont] at (0,-1.34) {$x^{[j]}$ 当前子步状态};
  \node[font=\fontsize{9.8pt}{11.7pt}\selectfont,text=SLBlue] at (-2.65,-1.81)
    {起始至当前坐标\quad 本轮新值};
  \node[font=\fontsize{9.8pt}{11.7pt}\selectfont,text=SLTextGray] at (2.65,-1.81)
    {后续至末位坐标\quad 前轮旧值};

  \node[sl634-card,minimum width=123mm,minimum height=11mm] at (0,-2.72) {};
  \node[inner sep=2.5pt,font=\fontsize{10.0pt}{12.0pt}\selectfont] (sweepend) at (-2.65,-2.87) {$x^{[d]}$};
  \node[inner sep=2.5pt,font=\fontsize{10.0pt}{12.0pt}\selectfont] (roundstate) at (.85,-2.87) {$x^{(t)}$};
  \node[inner sep=2.5pt,font=\fontsize{9.8pt}{11.7pt}\selectfont,text=SLTextGray] (sample) at (4.25,-2.87) {轮末样本};
  \draw[<->,draw=SLGray,line width=.60pt] (sweepend.east)--(roundstate.west)
    node[midway,above=4pt,font=\fontsize{9.6pt}{11.5pt}\selectfont,text=SLTextGray] {状态相同};
  \draw[-{Stealth[length=1.5mm,width=1.05mm]},draw=SLGray,line width=.60pt] (roundstate.east)--(sample.west)
    node[midway,above=4pt,font=\fontsize{9.6pt}{11.5pt}\selectfont,text=SLTextGray] {仅此记录};
\end{tikzpicture}
\captionsetup{width=.94\linewidth}
\caption[系统扫描逐坐标即时写回并区分子步状态与轮末样本]{系统扫描按固定次序即时写回；当前子步的前段使用本轮新值，后段沿用前轮旧值；末位更新结束后，末位状态与本轮样本状态相同并记录为轮末样本。}
\label{fig:V5-C04-coordinate-sweep}
\end{figure}

\FloatBarrier
\noindent\textbf{读图顺序。}先按上方编号与箭头确认$j=1,\ldots,d$的固定次序；再看$x^{[j]}$左侧坐标已写入同轮新值、右侧坐标仍保留上一轮旧值，因而后续更新会读取此前的新值；最终只把$x^{[d]}=x^{(t)}$记作本轮样本，不把中间状态当作一轮输出。\par
\paragraph{轮内状态与轮末样本。}中途停止得到的是部分扫描状态；只有完成第$d$个坐标更新后，$x^{[d]}=x^{(t)}$才能追加到轮末样本序列。
\end{document}
